\section{Introduction and statements of the results}
Let $V_n=\F^n$ with the standard bilinear form
$\langle x,y\rangle=\sum_{i=1}^n x_i y_i$.
The graph $\OG_n^*$ has the nonzero linear subspaces of $V_n$ as its vertices.
Distinct vertices $U,W$ are adjacent when every vector of $U$ is orthogonal
to every vector of $W$. Its induced subgraph on one-dimensional subspaces is
denoted by $\Gamma_n$. Over $\F$, each line has a unique nonzero vector.

Nikandish~\cite{Nikandish} asks for the clique number in Problem~4.2 and the
chromatic number for $n\geq6$ in Problem~4.3. Put
$N(r)=\sum_{d=1}^r {r\brack d}_2$, where ${r\brack d}_2$ is the Gaussian
binomial coefficient, and set $N(0)=0$. Our results are
\[
 \omega(\OG_n^*)=\max\{n,N(\lfloor n/2\rfloor)+(n\bmod2)\}\quad(n\geq1),
\]
and
\[
 \chi(\OG_6^*)=15,\qquad 12\leq\chi(\Gamma_6)\leq13.
\]
The first formula answers Problem~4.2 in all dimensions. The second answers
the dimension-six case of Problem~4.3 and quantifies the separation between
coloring the line graph and coloring the full subspace graph.

The clique argument is algebraic. The full-graph chromatic upper bound uses
an explicit finite certificate, while the line lower bound follows from the
geometry of pairwise nonorthogonal vectors. The verification details and the
scope of the Lean formalization are collected in Section~\ref{sec:verification}.
